\documentclass[aspectratio=169, 10pt]{beamer}

\usetheme{Madrid}
\usefonttheme{professionalfonts}

\usepackage[utf8]{inputenc}
\usepackage[indonesian]{babel}
\usepackage{amsmath,amsfonts,amssymb}
\usepackage{graphicx}
\usepackage{booktabs}
\usepackage{colortbl}
\usepackage{listings}
\usepackage{xcolor}

\definecolor{mainblue}{RGB}{24, 43, 73}
\definecolor{accentblue}{RGB}{41, 128, 185}
\definecolor{accentgold}{RGB}{230, 126, 34}
\definecolor{darkgray}{RGB}{44, 62, 80}
\definecolor{lightbg}{RGB}{248, 250, 252}
\definecolor{bordergray}{RGB}{200, 210, 220}
\definecolor{codebg}{RGB}{240, 243, 246}
\definecolor{codegreen}{RGB}{39, 174, 96}

\setbeamercolor{structure}{fg=mainblue}
\setbeamercolor{palette primary}{bg=mainblue, fg=white}
\setbeamercolor{palette secondary}{bg=accentblue, fg=white}
\setbeamercolor{palette tertiary}{bg=mainblue!80, fg=white}
\setbeamercolor{title}{fg=white, bg=mainblue}
\setbeamercolor{frametitle}{fg=white, bg=mainblue}
\setbeamercolor{block title}{bg=mainblue, fg=white}
\setbeamercolor{block body}{bg=lightbg, fg=darkgray}
\setbeamercolor{block title alerted}{bg=accentgold, fg=white}
\setbeamercolor{block body alerted}{bg=accentgold!10, fg=darkgray}

\lstset{
    language=Python,
    backgroundcolor=\color{codebg},
    basicstyle=\ttfamily\scriptsize\color{darkgray},
    keywordstyle=\bfseries\color{accentblue},
    stringstyle=\color{accentgold},
    commentstyle=\itshape\color{codegreen},
    showstringspaces=false,
    breaklines=true,
    frame=single,
    rulecolor=\color{bordergray},
    numbers=none
}

\title[\textbf{Deret Waktu (Minggu 11)}]{\textbf{Deep Learning Sekuensial: RNN, LSTM, dan GRU}}
\subtitle{Arsitektur Gerbang Memori, Masalah Vanishing Gradient, dan Peramalan Daya EBT Surya/Angin}
\author[\textbf{Ir. Novalio Daratha, Ph.D.}]{\textbf{Ir. Novalio Daratha, S.T., M.Sc., Ph.D.}}
\institute[FT UNIB]{Program Studi Teknik Elektro -- Fakultas Teknik\\Universitas Bengkulu}
\date{Minggu 11 / Semester Ganjil 2026}

\begin{document}

\begin{frame}[plain]
    \titlepage
\end{frame}

\begin{frame}{Agenda Perkuliahan Minggu ke-11}
    \begin{columns}[T]
        \begin{column}{0.48\textwidth}
            \begin{block}{Capaian Pembelajaran (Sub-CPMK 10)}
                Mahasiswa mampu menyusun tensor 3D sekuensial, memahami mekanisme gerbang arsitektur LSTM dan GRU, mengatasi masalah kelenyapan gradien (*vanishing gradient*), serta melatih jaringan *Deep Learning* untuk peramalan daya pembangkit EBT fluktuatif.
            \end{block}
        \end{column}
        \begin{column}{0.48\textwidth}
            \begin{enumerate}
                \item \textbf{Keterbatasan Jaringan Saraf Konvensional (MLP)}
                \item \textbf{Recurrent Neural Network (RNN) Standar}
                \item \textbf{Masalah Vanishing \& Exploding Gradient}
                \item \textbf{Arsitektur Long Short-Term Memory (LSTM)}
                \item \textbf{Arsitektur Gated Recurrent Unit (GRU)}
                \item \textbf{Implementasi PyTorch: Peramalan Daya PLTS}
            \end{enumerate}
        \end{column}
    \end{columns}
\end{frame}

\begin{frame}{Dari MLP ke Recurrent Neural Network (RNN)}
    \begin{columns}[T]
        \begin{column}{0.48\textwidth}
            \begin{block}{Mengapa MLP Kurang Efektif?}
                \begin{itemize}
                    \item Mengasumsikan seluruh sampel input independen satu sama lain.
                    \item Tidak memiliki memori internal untuk mempertahankan informasi riwayat urutan temporal yang panjang.
                \end{itemize}
            \end{block}
        \end{column}
        \begin{column}{0.48\textwidth}
            \begin{block}{Mekanisme RNN}
                Memiliki status tersembunyi (\textit{Hidden State} $h_t$) yang berulang (*recurrent loop*):
                \begin{equation}
                    h_t = \tanh(W_{hh} h_{t-1} + W_{xh} x_t + b_h)
                \end{equation}
                \begin{equation}
                    \hat{y}_t = W_{hy} h_t + b_y
                \end{equation}
            \end{block}
        \end{column}
    \end{columns}
\end{frame}

\begin{frame}{Arsitektur Long Short-Term Memory (LSTM)}
    Mengatasi *vanishing gradient* dengan jalur sel memori bebas hambatan (*Constant Error Carousel*):
    \begin{block}{Mekanisme 3 Gerbang LSTM}
        \begin{equation*}
        \begin{aligned}
            \text{Forget Gate: } &f_t = \sigma(W_f [h_{t-1}, x_t] + b_f) \quad \text{(Membuang memori usang)} \\
            \text{Input Gate: } &i_t = \sigma(W_i [h_{t-1}, x_t] + b_i), \quad \tilde{C}_t = \tanh(W_c [h_{t-1}, x_t] + b_c) \\
            \text{Cell State: } &C_t = f_t \odot C_{t-1} + i_t \odot \tilde{C}_t \quad \text{(Memperbarui akumulasi memori)} \\
            \text{Output Gate: } &o_t = \sigma(W_o [h_{t-1}, x_t] + b_o), \quad h_t = o_t \odot \tanh(C_t)
        \end{aligned}
        \end{equation*}
    \end{block}
\end{frame}

\begin{frame}{Gated Recurrent Unit (GRU): Varian Lebih Ringkas}
    Menggabungkan *Cell State* dan *Hidden State* serta mereduksi jumlah gerbang menjadi 2:
    \begin{columns}[T]
        \begin{column}{0.5\textwidth}
            \begin{block}{Persamaan GRU}
                \begin{equation*}
                \begin{aligned}
                    \text{Reset Gate: } &r_t = \sigma(W_r [h_{t-1}, x_t]) \\
                    \text{Update Gate: } &z_t = \sigma(W_z [h_{t-1}, x_t]) \\
                    \text{Candidate: } &\tilde{h}_t = \tanh(W [r_t \odot h_{t-1}, x_t]) \\
                    \text{Hidden State: } &h_t = (1 - z_t) \odot h_{t-1} + z_t \odot \tilde{h}_t
                \end{aligned}
                \end{equation*}
            \end{block}
        \end{column}
        \begin{column}{0.5\textwidth}
            \begin{alertblock}{Perbandingan Rekayasa: LSTM vs GRU}
                \begin{itemize}
                    \item \textbf{GRU}: Parameter 25\% lebih sedikit, konvergensi lebih cepat pada dataset terbatas (misal data sensor mikrogrid).
                    \item \textbf{LSTM}: Lebih tangguh menangani ketergantungan temporal jangka panjang yang sangat rumit.
                \end{itemize}
            \end{alertblock}
        \end{column}
    \end{columns}
\end{frame}

\begin{frame}[fragile]{Implementasi PyTorch: Arsitektur LSTM Model}
\begin{lstlisting}
import torch
import torch.nn as nn

class SolarLSTM(nn.Module):
    def __init__(self, n_features=4, hidden_size=64, num_layers=2, pred_len=24):
        super(SolarLSTM, self).__init__()
        self.lstm = nn.LSTM(input_size=n_features, hidden_size=hidden_size,
                            num_layers=num_layers, batch_first=True, dropout=0.2)
        self.fc = nn.Linear(hidden_size, pred_len)

    def forward(self, x):
        # Input x shape: (batch_size, seq_len, n_features)
        lstm_out, _ = self.lstm(x)
        # Ambil hidden state terakhir dari sequence
        last_hidden = lstm_out[:, -1, :]
        out = self.fc(last_hidden) # Output: (batch_size, pred_len)
        return out
\end{lstlisting}
\end{frame}

\begin{frame}{Rangkuman \& Persiapan Minggu ke-12}
    \begin{block}{Poin Kunci Minggu 11}
        \begin{itemize}
            \item LSTM dan GRU mengatasi kelemahan mendasar RNN klasik dalam menangkap memori jangka panjang.
            \item Gerbang non-linier bertindak sebagai saklar kendali aliran informasi internal jaringan.
            \item Penskalaan data 3D Tensor $(N, T, D)$ wajib dinormalisasi secara hati-hati agar gradien stabil.
        \end{itemize}
    \end{block}
    \vspace{0.2cm}
    \begin{alertblock}{Materi Minggu Depan (Minggu ke-12)}
        \textbf{Arsitektur Sekuensial Lanjut}: 1D-Temporal CNN, Encoder-Decoder (Seq2Seq), dan Attention Mechanism / Transformer untuk peramalan multi-bus simultan.
    \end{alertblock}
\end{frame}

\begin{frame}[plain]
    \begin{center}
        \vspace{1.5cm}
        {\Huge \color{mainblue} \textbf{Terima Kasih}}\\[0.4cm]
        {\large Pertanyaan \& Diskusi Deep Learning Sekuensial}\\[1.2cm]
        \textbf{Laboratorium Sistem Tenaga \& Komputasi Cerdas}\\
        Jurusan Teknik Elektro -- Universitas Bengkulu
    \end{center}
\end{frame}

\end{document}
